% DSAA2011 Course Project (Task A) -- report draft
% ---------------------------------------------------------------------------------
% FORMAT NOTE: the course requires its own LaTeX style file (link in the project PDF).
% Put that .sty file next to this .tex and replace COURSE_STYLE below by its name
% (e.g. neurips_2024). Do NOT use the "preprint" option. Until the file is present,
% the \IfFileExists fallback compiles with plain article + 1-inch margins.
% ---------------------------------------------------------------------------------
\documentclass{article}
\IfFileExists{COURSE_STYLE.sty}{\usepackage[final]{COURSE_STYLE}}{\usepackage[margin=1in]{geometry}}
\usepackage[utf8]{inputenc}
\usepackage[T1]{fontenc}
\usepackage{lmodern}
\usepackage{hyperref}
\usepackage{url}
\usepackage{booktabs}
\usepackage{amsmath}
\usepackage{microtype}
\usepackage{xcolor}
\usepackage{graphicx}
\usepackage{caption}
\usepackage{subcaption}
\usepackage[numbers]{natbib}
\graphicspath{{../figures/}}

\title{``Where am I?'' Identifying Australia's Climate Zone\\from Daily Weather Observations}
\author{Group XX --- DSAA2011 (L02), HKUST(GZ)\\
  \texttt{Member A, Member B, Member C, Member D}}

\begin{document}
\maketitle

\begin{abstract}
We turn the \emph{Rain in Australia} dataset (145,460 daily observations at 49 stations, 2007--2017) into a
five-class problem: given one day of weather at a station, predict the station's climate zone, defined as five
project-specific groups of K\"oppen-Geiger codes (Tropical, Arid, Subtropical, Mediterranean, Oceanic). The labels
are derived reproducibly from Bureau of Meteorology coordinates and a 100 arc-second K\"oppen raster; the full
mapping is listed in Appendix~\ref{app:map}. Because the zone is a property of the station, we pay special attention
to \emph{station leakage}: missing-value indicators identify stations, and every model is evaluated on a random
split, a time split and on stations never seen in training. t-SNE and clustering show that single days group by
weather regime rather than by climate; only the tropics separate clearly. Over 15 groups of unseen stations,
logistic regression and a tuned decision tree reach a macro-F1 of $0.51\pm0.06$ and $0.53\pm0.06$ (0.62--0.69 on a
random split), gradient boosting $0.60\pm0.07$, and gradient boosting on 14-day station means (``climate memory'')
$0.68\pm0.09$; both improvements hold in every one of the 15 folds.
\end{abstract}

% =================================================================================
\section{Introduction}
Most work on this dataset \citep{kaggle} predicts \texttt{RainTomorrow}. We ask a different question that connects all parts of
the project: \emph{does a day of weather carry the signature of the climate it comes from?} K\"oppen zones are
defined by long-term monthly statistics, so the answer is not obvious, and it tests how far simple models can go
from a single observation.

\textbf{Main results.}
(1) The missing-value pattern alone identifies the station ten times better than chance, so whole-station sensor
columns and missing indicators are excluded.
(2) t-SNE separates only the Tropical zone (96\% of tropical days have a tropical majority among their neighbours);
the visible seasonal gradient comes mostly from the day-of-year inputs.
(3) The best clustering of days (K-means, $k=3$) finds weather regimes, not climate zones (ARI $\le 0.13$);
clustering \emph{station climate profiles} recovers the tropics and the subtropical coast exactly (ARI 0.35).
(4) A fully grown decision tree overfits (train accuracy 1.000, test 0.652) and its advantage over logistic
regression disappears on unseen stations; validation fixes most of the overfitting.
(5) Flexible models gain much less on unseen stations than on a random split; a time split shows that this
optimism comes mostly from knowing the station, not from adjacent days. 14-day trailing means raise the
unseen-station macro-F1 of gradient boosting by $0.083\pm0.039$ (better in 15 of 15 folds).
All code is in \texttt{project\_XX\_TaskA.ipynb}; every number below is printed there.

% =================================================================================
\section{Mandatory tasks}

\subsection{Target construction and data preprocessing}
\textbf{Target.} We matched each \texttt{Location} by hand to a Bureau of Meteorology (BoM) station
\citep{bom_stations} and looked up the K\"oppen-Geiger code of the pixel nearest to its coordinates with
\texttt{kgcpy} 1.1.8, the Python port of the R package \texttt{kgc} \citep{kgc}. Both ship a global $12960\times6480$
raster (100 arc-seconds, $\approx$3\,km per pixel) and cite Rubel et al.~\citep{rubel2017} as its source; that paper is
titled after the European Alps but describes the downscaling used for the authors' re-analysed global map, which
represents 1986--2010 \citep{kgmaps}. Our five classes are \emph{project-defined groups} of K\"oppen codes
\citep{kottek2006}: groups A and B are kept whole, and group C is split by K\"oppen's own second and third letters
into Cfa (Subtropical), Cs (Mediterranean) and Cfb (Oceanic). Only Am, Aw, BWh, BSh, BSk, Cfa, Csa, Csb and Cfb occur;
there are no D or E climates, and any other code would receive no label. Appendix Table~\ref{tab:map} lists every
station with its BoM number, coordinates, code and group. Norfolk Island receives the code ``Ocean'': the island is
absent from the raster, and the nearest land pixel lies 736\,km away in northern New Zealand, so this is not a
coordinate error. It is dropped (3,009 rows; it had no label, so the class shares of the rest are unchanged), leaving
142,451 rows from 48 stations. Classes range from 7.6\% (Tropical, 4 stations) to 32.2\% (Subtropical, 15 stations);
the majority baseline accuracy is 0.322.

\textbf{Leakage audit.} \texttt{Location} (the zone is a function of it) and \texttt{RainTomorrow} (future) are
removed; \texttt{RISK\_MM} is checked for and absent in Kaggle version 2. \texttt{Location} is kept outside the
feature table and used only for grouped splits.

\textbf{Missing values.} Missingness is a station property (Fig.~\ref{fig:missing}): Sunshine, Evaporation and
Cloud are missing for 12--19 entire stations. A decision tree trained on the 20 missing flags \emph{only} predicts
the zone with accuracy 0.439 (baseline 0.322) and the exact station with 0.240 (chance 0.024). Adding missing
indicators to logistic regression raises random-split accuracy from 0.606 to 0.671 but leaves unseen-station
accuracy unchanged (0.569 $\rightarrow$ 0.572), i.e.\ the gain is pure station recognition. We therefore drop the
four sensor columns and use median imputation without indicators. Two ``missing'' values are not missing: every
day with wind speed 0 has no wind direction (calm), and \texttt{RainToday} equals \texttt{Rainfall} $>1$\,mm on
100\% of rows (dropped as redundant).

\textbf{Encoding and scaling.} Wind direction and speed are combined into vectors
$(u,v)=(-s\sin\theta,-s\cos\theta)$, with calm $=(0,0)$; this keeps the circular order of the 16 compass points with
6 numbers instead of 48 one-hot columns. Day of year is encoded as $(\sin,\cos)$; the year is not used.
Rainfall is $\log(1+x)$-transformed (skew 9.9 $\rightarrow$ 2.1). Median imputation and standardisation sit inside
scikit-learn pipelines \citep{sklearn} fitted on training data only. The final \texttt{core} set has 20 features.

\textbf{Observed patterns.} The seasonal cycles follow the K\"oppen definitions (Fig.~\ref{fig:season}): summer
rain in the tropics, winter rain in the Mediterranean zone, the lowest humidity in the Arid zone, the coolest
summers in the Oceanic zone. Temperature variables are strongly correlated ($r=0.985$ for MaxTemp--Temp3pm).

\begin{figure}[t]
  \centering
  \begin{subfigure}{0.44\linewidth}\includegraphics[width=\linewidth]{p1_missing_heatmap}
    \caption{Missing rate per station (rows, coloured by zone).}\label{fig:missing}\end{subfigure}\hfill
  \begin{subfigure}{0.54\linewidth}\includegraphics[width=\linewidth]{p1_seasonal_cycle}
    \caption{Monthly means by zone.}\label{fig:season}\end{subfigure}
  \caption{Preprocessing insights: missingness is a station fingerprint; the seasonal cycle separates zones.}
\end{figure}

\subsection{Data visualization with t-SNE}
We embedded a balanced sample of 6,000 days (1,200 per zone) with t-SNE \citep{tsne} (perplexity 30, PCA
initialisation). To avoid over-reading the picture we report \emph{k-NN purity} (5-fold accuracy of a 10-NN
classifier on the coordinates), the same predictions scored with Cohen's $\kappa$ (which corrects for the different
chance levels of zone, season and station labels) and trustworthiness \citep{venna2001}. These are diagnostics of
the embedding, not estimates of generalisation: t-SNE used all points at once and days of one station fall on both
sides of every fold.

Only the Tropical zone forms its own region (Fig.~\ref{fig:tsne}): 96\% of tropical days have a tropical majority
among their 10 neighbours, against 0.50--0.64 for the other zones. The month colouring shows a summer--winter
gradient, but with all features neighbourhoods agree equally on season and zone ($\kappa$ 0.55 vs 0.54), and without
the two day-of-year inputs season agreement falls to $\kappa=0.17$ while zone agreement stays at 0.49: the weather
itself says more about the zone than about the season. Summer-only and winter-only embeddings barely change zone
$\kappa$ (0.57, 0.55) but show which zone is distinct when: tropical days form an island in winter, arid days move
to their own edge in summer. The most mixed pair is Mediterranean--Oceanic. Station purity is 0.31 ($\kappa$ 0.29;
chance 0.02), so station-specific information is present in the weather variables---a warning, tested in
Sections~2.4--3, that random splits could reward recognising stations. PCA keeps 42\% of the variance in 2-D and
preserves neighbourhoods far worse than t-SNE (trustworthiness 0.809 vs 0.985; zone $\kappa$ 0.22 vs 0.54); this
shows that t-SNE keeps local neighbourhoods better, not by itself that the structure is non-linear. Purity varies by
about $\pm0.01$ across perplexities 5--100 and three seeds.

\begin{figure}[t]
  \centering
  \includegraphics[width=\linewidth]{p2_tsne_main}
  \caption{t-SNE of 6,000 days coloured by zone, month and maximum temperature.}\label{fig:tsne}
\end{figure}

\subsection{Clustering analysis}
\textbf{Algorithms.} K-means (centroid model), Ward agglomerative clustering (deterministic, hierarchical, same
variance criterion as K-means) and a Gaussian mixture (GMM; full covariance, \texttt{reg\_covar} $=10^{-6}$,
2 initialisations; elliptical clusters, BIC). Each was run for $k=2,\dots,10$ on the same 6,000 days and evaluated
with silhouette \citep{silhouette}, Davies--Bouldin \citep{daviesbouldin}, Calinski--Harabasz, inertia/BIC, ARI
\citep{ari} and NMI against zone and season. \emph{Stability} is measured identically for all three algorithms: refit
on ten 80\% subsamples and compute the ARI with the reference partition on the shared days. All curves are in
Appendix Fig.~\ref{fig:ksel}.

\textbf{Selection rule and result.} Among solutions with stability $\ge0.90$ we choose the best mean rank on
silhouette, Davies--Bouldin and Calinski--Harabasz; external metrics are not used. (The rule was written down
after our first analysis and reproduces its choice.) Eligible are K-means $k=2$--4 and GMM $k=3,5$; Ward is never
eligible (0.26--0.36 for every $k$: removing 20\% of the days changes its greedy early merges). The rule selects
\textbf{K-means with $k=3$} (mean rank 1.67): the lowest Davies--Bouldin of all 27 solutions (1.88), silhouette 0.164
(0.171 at $k=2$, which ranks only sixth on Davies--Bouldin) and stability 0.98 (Table~\ref{tab:clust}). The GMM is
degenerate: one component (weight 0.615) contains exactly the dry days and has zero rainfall variance, because 63\% of
days have Rainfall $=0$; this point mass explains why BIC keeps improving while silhouette falls. The GMM fits at $k=3$ and 5
converged (24--40 EM iterations); raising \texttt{reg\_covar} to 0.1 removes the zero variance and triples the
silhouette at $k=3$ (0.032 $\rightarrow$ 0.100), but GMM remains clearly behind K-means.

\textbf{Interpretation.} The three clusters are weather regimes (Fig.~\ref{fig:clprof}): warm days (40\% tropical),
cool calm high-pressure days (mostly winter/autumn) and cool, windy, wet frontal days (mostly Oceanic and
Mediterranean). Every zone has all three kinds of days, so ARI with the zone never exceeds 0.13 for any algorithm or
$k$. As an extra analysis we clustered the 48 \emph{stations}, each described by 12 monthly means of four variables.
Ward with $k=5$ recovers all 4 tropical stations and the 9 coastal subtropical stations exactly
(Fig.~\ref{fig:dendro}); the remaining clusters follow coast versus inland rather than K\"oppen's rainfall-season
rule. ARI rises from about 0.1 (days) to 0.345 (stations): station profiles \emph{partly} recover the climate
structure, consistent with climate being a statistic of many days.

\begin{table}[t]
  \centering\small
  \caption{Clustering of 6,000 days at $k=3$, and station-profile clustering ($k=5$).}\label{tab:clust}
  \setlength{\tabcolsep}{4pt}\footnotesize
  \begin{tabular}{lcccccc}
    \toprule
    & Silhouette $\uparrow$ & Davies--Bouldin $\downarrow$ & Calinski--H. $\uparrow$ & ARI zone & NMI zone & Stability \\
    \midrule
    K-means ($k=3$) & \textbf{0.164} & \textbf{1.880} & \textbf{1086} & 0.098 & 0.122 & \textbf{0.983} \\
    Ward ($k=3$)    & 0.117 & 2.150 & 722  & 0.114 & 0.150 & 0.358 \\
    GMM ($k=3$; BIC 181,637) & 0.032 & 2.958 & 347  & 0.034 & 0.097 & 0.904 \\
    \midrule
    Ward, 48 station profiles ($k=5$) & 0.319 & -- & -- & \textbf{0.345} & \textbf{0.540} & -- \\
    \bottomrule
  \end{tabular}
\end{table}

\begin{figure}[t]
  \centering
  \includegraphics[width=\linewidth]{p3_cluster_profiles}
  \caption{K-means ($k=3$): cluster profiles (z-score colour, original units as numbers), zone mix and season mix.}\label{fig:clprof}
\end{figure}
\begin{figure}[t]
  \centering
  \includegraphics[width=0.92\linewidth]{p3_station_dendrogram}
  \caption{Ward dendrogram of 48 station climate profiles; labels coloured by K\"oppen zone.}\label{fig:dendro}
\end{figure}

\subsection{Prediction: training and testing}
\textbf{Target and models.} Target \texttt{Zone}; features: the 20 core features. We chose multinomial
\textbf{logistic regression} (linear baseline, probabilistic output, interpretable weights) and a \textbf{CART decision
tree} (non-linear, readable rules), first with default settings ($C=1$; unlimited depth).

\textbf{Splits.} The required stratified random 70/30 split answers ``can the model classify new days of known
stations?''. A station-grouped 70/30 split (\texttt{GroupShuffleSplit}) answers ``can it classify stations it has
never seen?''. Its 15 test stations are Cairns (Tropical, 3,040 days); Alice Springs, Nhil, Uluru (Arid, 6,196);
Badgerys Creek, Newcastle, Norah Head, Richmond, Sydney Airport (Subtropical, 15,070); Dartmoor, Nuriootpa
(Mediterranean, 6,018); Bendigo, Melbourne Airport, Tuggeranong, Watsonia (Oceanic, 12,097). Because one draw of
15 stations is noisy, Section~3 repeats the station-grouped evaluation over 15 folds.

\textbf{Training and results.} Each model is one pipeline \texttt{fit} on the training rows. The tree grows to depth
35 with 18,947 leaves (about 5 training days per leaf). Table~\ref{tab:p4} gives accuracy and macro-F1 on the
training, test and entire sets; Fig.~\ref{fig:cm} shows the six confusion matrices of the random split.

\begin{table}[t]
  \centering\small
  \caption{Default models on the training, test and entire sets (accuracy / macro-F1).}\label{tab:p4}
  \begin{tabular}{llccc}
    \toprule
    Split & Model & Train & Test & Entire \\
    \midrule
    Random  & LR & 0.604 / 0.620 & 0.599 / 0.615 & 0.602 / 0.618 \\
    Random  & DT & 1.000 / 1.000 & 0.652 / 0.669 & 0.895 / 0.900 \\
    Grouped & LR & 0.609 / 0.625 & 0.563 / 0.561 & 0.596 / 0.609 \\
    Grouped & DT & 1.000 / 1.000 & 0.533 / 0.536 & 0.861 / 0.865 \\
    \bottomrule
  \end{tabular}
\end{table}

\begin{figure}[t]
  \centering
  \includegraphics[width=0.66\linewidth]{p4_confusion_random}
  \caption{Confusion matrices (row-normalised, counts in brackets) for the random split: training, test and entire set.}\label{fig:cm}
\end{figure}

\textbf{Interpretation.} LR has nearly identical scores on all three sets: with 105 parameters it cannot memorise
100,000 days; its errors are bias. The default tree is perfect on the training set but 0.35 lower on the test set,
a clear case of overfitting, yet still better than LR on known stations. The entire set is
$0.7\cdot$train$+0.3\cdot$test and is not an independent estimate (0.895 for the tree is misleading). On unseen
stations the tree falls below LR (0.533 vs 0.563): part of what it learned was station-specific. Tropical is
recognised best (recall 0.77--0.92), its errors going to Subtropical; the largest confusion is Mediterranean
$\leftrightarrow$ Oceanic (up to 41\% of Mediterranean test days on unseen stations); Arid has the lowest recall and
is mostly confused with the inland Subtropical stations.

\textbf{Decision boundary.} Fig.~\ref{fig:boundary} retrains each model on the two most informative interpretable
features, MinTemp and Humidity3pm. LR draws straight boundaries and assigns no region to Mediterranean at all;
the default tree fragments the plane into tiny rectangles; a depth-6 tree keeps the large structures, carves out a
Mediterranean box and scores best (two-feature test macro-F1 0.474 vs 0.412 and 0.390).

\begin{figure}[t]
  \centering
  \includegraphics[width=0.88\linewidth]{p4_decision_boundary}
  \caption{Decision boundaries on MinTemp $\times$ Humidity3pm (1,500 random test days).}\label{fig:boundary}
\end{figure}

\subsection{Evaluation and choice of prediction model}
\textbf{Metrics.} Per-zone precision, recall and F1 were computed directly from the confusion matrices
($\mathrm{TP}_k=C_{kk}$, column sums for precision, row sums for recall) and checked against scikit-learn. On the
random test set LR obtains macro-F1 0.615 and weighted-F1 0.593, the tree 0.669 and 0.652. Macro-F1 exceeds
weighted-F1 because the smallest class, Tropical, is the easiest. Arid and Mediterranean have the lowest F1.

\textbf{ROC and AUC.} We use one-vs-rest ROC curves (one per zone) and macro-AUC. The default tree has a lower
macro-AUC than LR (0.788 vs 0.859) despite higher accuracy: its pure leaves give 0/1 probabilities and therefore
almost no ranking.

\textbf{Validation.} We tuned LR ($C$, class weights) and the tree (max depth, min samples per leaf) with 5-fold
grid search on macro-F1, once with stratified folds on the random training set and once with
\texttt{StratifiedGroupKFold} by station on the grouped training set (Fig.~\ref{fig:val}). With one sample per leaf
the tree's validation score peaks at depth 13--16 and then declines while training F1 reaches 1.0; both CV schemes
select depth 16 with at least 20 days per leaf, which keeps the curve at its peak. The grouped validation level is
0.26 lower and much noisier. LR is flat for $C\ge 0.01$; under grouped CV even LR has a train--validation gap of
about 0.16, which reflects distribution shift between stations rather than overfitting.

\begin{figure}[t]
  \centering
  \includegraphics[width=0.85\linewidth]{p5_validation_curves}
  \caption{Validation curves for the decision tree (left) and logistic regression (right), random vs.\ station-grouped CV.}\label{fig:val}
\end{figure}

\begin{table}[t]
  \centering\small
  \caption{Default vs.\ tuned models on the held-out test sets.}\label{tab:p5}
  \begin{tabular}{lllcccc}
    \toprule
    Split & Model & Version & Train mF1 & Test mF1 & Gap & Test macro-AUC \\
    \midrule
    Random  & LR & default / tuned & 0.620 / 0.620 & 0.615 / 0.615 & 0.005 / 0.005 & 0.859 / 0.859 \\
    Random  & DT & default / tuned & 1.000 / 0.758 & 0.669 / \textbf{0.691} & 0.331 / 0.067 & 0.788 / \textbf{0.889} \\
    Grouped & LR & default / tuned & 0.625 / 0.615 & 0.561 / 0.554 & 0.064 / 0.062 & 0.823 / 0.821 \\
    Grouped & DT & default / tuned & 1.000 / 0.778 & 0.536 / 0.559 & 0.464 / 0.219 & 0.705 / 0.813 \\
    \bottomrule
  \end{tabular}
\end{table}

\begin{figure}[t]
  \centering
  \includegraphics[width=0.55\linewidth]{p5_roc_tuned}
  \caption{One-vs-rest ROC curves of the tuned models: random test set (top) and 15 unseen stations (bottom).}\label{fig:roc}
\end{figure}

\textbf{Discussion.} Both simple models recognise tropical days reliably (one-vs-rest AUC above 0.97) but confuse
the three southern zones, which overlap in feature space as already seen with t-SNE. Logistic regression shows
almost no gap between training and test scores: it underfits because the zone boundaries are not linear, and
tuning $C$ does not help. The unconstrained decision tree is the opposite: perfect on the training set, clearly
worse on the test set, and its near-binary probabilities give a poor ROC curve. Limiting depth and leaf size through
cross-validation removes most of this overfitting and raises both macro-F1 and AUC (Table~\ref{tab:p5},
Fig.~\ref{fig:roc}). The most important result is the comparison of the two splits. On a random split every station
is already known, so the tree partly memorises station-specific weather. On stations never seen in training, all
scores drop and vary strongly between stations, and the tuned tree is no longer clearly better than logistic
regression. Accuracy alone would hide these differences; macro-F1 and per-zone AUC reveal them. Further improvement
needs more flexible models and features that summarise more than one day of weather.

\textbf{Pooled vs.\ per-fold scores.} Out-of-fold predictions from 5-fold \texttt{StratifiedGroupKFold} over all 48
stations can be summarised by one pooled macro-F1 (LR 0.553) or per fold (LR $0.544\pm0.080$, range 0.44--0.63). Per
fold, macro-F1 must be averaged over the classes \emph{present}: two folds contain no tropical station, and
scikit-learn's default then counts Tropical with $F_1=0$ whenever it is predicted, which lowers the mean to 0.483.
This bias also explains the low grouped-CV levels of the grid searches; it affects all candidates alike, so it does not
change their ranking.

\textbf{Error analysis.} The same out-of-fold predictions (each station predicted by a model that never saw it)
give accuracy 0.53 (LR) and 0.55 (tree), with per-station accuracy from below 0.1 to above 0.9. The worst stations
are the inland Subtropical stations (Wagga Wagga, Moree, Albury), Mediterranean stations of south-east South
Australia and western Victoria (predicted Oceanic) and cool semi-arid stations (Nhil, Salmon Gums) --- the same
coast/inland grouping found by station clustering.

% =================================================================================
\section{Open-ended exploration}
All experiments use the same 20 core features, the same median imputation fitted inside each training fold, the
same random 70/30 split and the same station folds; they report macro-F1 on the random test set and out-of-fold
macro-F1 over all 48 stations with station-grouped 5-fold CV (``unseen''). Model configurations, preprocessing and fit
times are listed in Appendix Table~\ref{tab:hyper}.

\textbf{Model comparison and complexity.} Fig.~\ref{fig:zoo} compares seven model classes. Flexible models (gradient
boosting \citep{friedman2001,lightgbm}, random forest \citep{breiman2001}, MLP, $k$-NN, LR with degree-2 polynomial
features) beat plain LR by 0.14--0.18 on the random split but only by 0.06--0.11 on unseen stations; the gap between
the two protocols grows with flexibility (LR 0.09, $k$-NN 0.18). Polynomial features with a stronger L2 penalty
($C=0.1$) remove much of LR's bias, and this gain transfers to unseen stations (0.525 $\rightarrow$ 0.623).
Histogram gradient boosting (HGB) is best under both protocols (0.795 / 0.638).

\textbf{Hyperparameter tuning.} A randomised search over 16 HGB configurations with station-grouped CV improved
validation macro-F1 by only 0.009 (test +0.004): all configurations fit the training stations to 0.93--0.97 but score
0.47--0.48 on validation stations, so the limit is station diversity, not hyperparameters.

\textbf{Feature engineering: climate memory.} Four derived features (diurnal range, humidity drop, pressure and
temperature change) change little (+0.003). We then replace each of the 18 daily features of station $s$ on day $t$
by its trailing calendar-day mean
\begin{equation*}
\bar x^{(N)}_{s,t,j}=\frac{1}{|W|}\sum_{\tau\in W}x_{s,\tau,j},\qquad W=\{\tau: t-N<\tau\le t,\ x_{s,\tau,j}\ \text{observed}\},
\end{equation*}
defined if $|W|\ge\max(1,\lfloor N/3\rfloor)$ and missing otherwise (then median-imputed inside the training fold);
day-of-year is kept. The window contains day $t$ and the $N-1$ days before it, never a later day; gaps in the record
shorten $|W|$ instead of extending the window. Computing the windows once on the full table does not leak under
station-grouped CV: a row's value depends only on its own station's weather, never on labels or other stations, and a
station is entirely in one fold. The experiment asks how many days of observation are needed to recognise the
climate (Fig.~\ref{fig:memory}). On unseen stations macro-F1 rises from 0.638 (1 day) to 0.710 (7 days) and 0.727
(14 days), then flattens; a 365-day mean drops to 0.655 because it removes the seasonal cycle on which the K\"oppen
definitions rely. Under a random split the same features reach 0.965 (14 days) and 0.99 (30 days or more), because
overlapping windows give every test day a near-copy in training---the strongest leakage example in the project.

\textbf{Feature selection.} Permutation importance of a multi-scale HGB (today + 7/30/90-day means, 74 features) on
unseen stations ranks 90-day means of wind, humidity, rainfall and minimum temperature highest; 20 selected features
match all 74 (0.718). Wind may partly reflect anemometer siting, but because importance is measured on unseen stations
it also reflects real wind regimes (trade winds, westerlies).

\textbf{Three evaluation protocols.} Table~\ref{tab:proto} and Fig.~\ref{fig:proto} evaluate the four key models
with (i) the random split, (ii) a time split (train before 2015-01-01, test from 2015-01-15 with a 14-day embargo, all
48 stations on both sides) and (iii) station-grouped 5-fold CV repeated with three fold seeds (15 held-out station
groups). Going from random to time costs only 0.016--0.033 for single-day models (0.07 for 14-day means); going from
time to unseen stations costs a further 0.09--0.21. The optimism of random splits therefore comes mainly from knowing
the station, not from adjacent days. The seed-42 folds used above were slightly favourable (0.01--0.04). Station
groups differ greatly in difficulty but affect all models alike, so paired differences are precise: HGB beats LR and
the tuned tree in 15 of 15 folds (vs.\ LR $+0.091\pm0.025$), and 14-day means beat single days in 15 of 15
($+0.083\pm0.039$), whereas the tree beats LR in only 10 of 15 ($+0.019\pm0.028$).

\begin{table}[t]
  \centering\small
  \caption{Macro-F1 of the key models under three protocols. Grouped: per-fold mean $\pm$ std over 15 held-out station
  groups (3 seeds $\times$ 5 folds), classes present in each fold.}\label{tab:proto}
  \begin{tabular}{lccc}
    \toprule
    Model & Random 70/30 & Time 70/30 (embargo 14 d) & Unseen stations \\
    \midrule
    Logistic regression ($C=1$)           & 0.615 & 0.599 & $0.509\pm0.058$ \\
    Decision tree (depth 16, leaf 20)     & 0.691 & 0.658 & $0.528\pm0.056$ \\
    HGB, daily features                   & 0.795 & 0.765 & $0.599\pm0.069$ \\
    HGB, 14-day trailing means            & 0.965 & 0.895 & $\mathbf{0.683\pm0.091}$ \\
    \bottomrule
  \end{tabular}
\end{table}

\begin{figure}[t]
  \centering
  \includegraphics[width=0.62\linewidth]{p6_protocols}
  \caption{Unseen-station macro-F1 of 15 held-out station groups (boxes, dots) and the two known-station protocols.}\label{fig:proto}
\end{figure}

\textbf{Final model.} For a single day we choose HGB (0.795 random, 0.765 time split, $0.60\pm0.07$ unseen stations);
with two weeks of observations, HGB on 14-day means ($0.68\pm0.09$ on unseen stations).

\begin{figure}[t]
  \centering
  \begin{subfigure}{0.53\linewidth}\includegraphics[width=\linewidth]{p6_model_comparison}
    \caption{Seven model classes, random vs.\ unseen stations.}\label{fig:zoo}\end{subfigure}\hfill
  \begin{subfigure}{0.45\linewidth}\includegraphics[width=\linewidth]{p6_climate_memory}
    \caption{Macro-F1 vs.\ number of days averaged.}\label{fig:memory}\end{subfigure}
  \caption{Open-ended exploration.}
\end{figure}

% =================================================================================
\section{Conclusion}
A single day of Australian weather carries a weak but real climate signature. Tropical days are unmistakable; the
three southern zones overlap because each of them has warm, calm and stormy days, and the K\"oppen boundaries between
them are thresholds on monthly statistics. Every part of the project points to the same distinction between
\emph{weather} and \emph{climate}: t-SNE and day-level clustering find seasons and weather regimes, station-level
clustering finds climate regions, and classifiers improve most when they see the average of a few weeks. The second
lesson is methodological: because the label belongs to the station, a random split rewards models for recognising
stations. Missing indicators, deep trees, $k$-NN and rolling windows all look excellent on a random split and much
less so on unseen stations, which is the setting that answers our question; a time split shows that adjacent days
account for only a small part of this gap. On 15 groups of unseen stations the best single-day model reaches a
macro-F1 of $0.60\pm0.07$ and two weeks of observations $0.68\pm0.09$. The main limitations are the small number of
stations (four tropical, seven arid), which makes every unseen-station estimate noisy, and the labels themselves: they
come from one K\"oppen raster whose boundaries carry their own uncertainty, which we did not test against an
independent map. More stations, or the BoM's full network, would be the most effective way to improve
generalisation.

% =================================================================================
\begin{thebibliography}{99}\small
\bibitem{kaggle} Kaggle user \texttt{jsphyg}. Rain in Australia (weatherAUS; observations from the Australian Bureau of Meteorology, distributed with the R package \texttt{rattle}), dataset \texttt{jsphyg/weather-dataset-rattle-package}, version 2. \url{https://www.kaggle.com/datasets/jsphyg/weather-dataset-rattle-package}.
\bibitem{bom_stations} Bureau of Meteorology. Station list, product IDCJMC0014. \url{ftp://ftp.bom.gov.au/anon2/home/ncc/metadata/sitelists/stations.zip}, accessed 2026-11-09.
\bibitem{rubel2017} F. Rubel, K. Brugger, K. Haslinger, I. Auer. The climate of the European Alps: shift of very high resolution K\"oppen-Geiger climate zones 1800--2100. \emph{Meteorologische Zeitschrift}, 26(2):115--125, 2017. doi:10.1127/metz/2016/0816.
\bibitem{kgc} C. Bryant, N. R. Wheeler, F. Rubel, R. H. French. kgc: Koeppen-Geiger Climatic Zones. R package version 1.0.0.2, 2017. doi:10.32614/CRAN.package.kgc. Python port \texttt{kgcpy} 1.1.8 (X. Yu, C. Bryant, N. R. Wheeler, F. Rubel, J. Ascencio Vasquez, R. H. French), \url{https://github.com/cwru-sdle/kgcpy}.
\bibitem{kgmaps} F. Rubel et al. World Maps of K\"oppen-Geiger climate classification: high resolution map and data (version March 2017), period 1986--2010. \url{https://koeppen-geiger.vu-wien.ac.at/present.htm}, accessed 2026-11-16.
\bibitem{kottek2006} M. Kottek, J. Grieser, C. Beck, B. Rudolf, F. Rubel. World map of the K\"oppen-Geiger climate classification updated. \emph{Meteorologische Zeitschrift}, 15(3):259--263, 2006.
\bibitem{sklearn} F. Pedregosa et al. Scikit-learn: machine learning in Python. \emph{JMLR}, 12:2825--2830, 2011.
\bibitem{tsne} L. van der Maaten, G. Hinton. Visualizing data using t-SNE. \emph{JMLR}, 9:2579--2605, 2008.
\bibitem{venna2001} J. Venna, S. Kaski. Neighborhood preservation in nonlinear projection methods: an experimental study. \emph{ICANN}, 2001.
\bibitem{silhouette} P. J. Rousseeuw. Silhouettes: a graphical aid to the interpretation and validation of cluster analysis. \emph{J. Comput. Appl. Math.}, 20:53--65, 1987.
\bibitem{daviesbouldin} D. L. Davies, D. W. Bouldin. A cluster separation measure. \emph{IEEE TPAMI}, 1(2):224--227, 1979.
\bibitem{ari} L. Hubert, P. Arabie. Comparing partitions. \emph{Journal of Classification}, 2:193--218, 1985.
\bibitem{breiman2001} L. Breiman. Random forests. \emph{Machine Learning}, 45:5--32, 2001.
\bibitem{friedman2001} J. H. Friedman. Greedy function approximation: a gradient boosting machine. \emph{Annals of Statistics}, 29(5):1189--1232, 2001.
\bibitem{lightgbm} G. Ke et al. LightGBM: a highly efficient gradient boosting decision tree. \emph{NeurIPS}, 2017.
\end{thebibliography}

% =================================================================================
\section*{Credit}
\textbf{Group members} (to be completed by the group): Member A --- \dots\,\%; Member B --- \dots\,\%;
Member C --- \dots\,\%; Member D --- \dots\,\%.

\textbf{Use of GenAI tools.} The teaching team allowed free use of AI tools provided that a clear log is kept.
We used \emph{Claude (Anthropic)} to analyse the project brief, brainstorm design options, write and debug the
notebook code, generate figures and draft this report. Every session is recorded in \texttt{AI\_LOG.md} (date, purpose,
output, our changes and verification). All reported numbers were checked against the executed notebook; design
decisions and their reasons are recorded in \texttt{DECISIONS.md}. \emph{(Group: state the approximate share of
AI-generated code and text here.)}

% =================================================================================
\appendix
\section{Appendix}\label{app:map}

\begin{table}[!htbp]
  \centering\scriptsize\setlength{\tabcolsep}{2.5pt}
  \caption{Label audit: dataset \texttt{Location} $\rightarrow$ BoM site number $\rightarrow$ coordinates (BoM station
  list) $\rightarrow$ K\"oppen code of the nearest pixel $\rightarrow$ project group. $^\ast$: a neighbouring land pixel
  (3$\times$3 window) belongs to a different group. Subtrop.\ = Subtropical, Medit.\ = Mediterranean. Generated by the notebook
  (Section 0.2) and stored in \texttt{data/processed/location\_zone.csv}.}\label{tab:map}
\begin{tabular}{llrrll}\toprule Location & BoM & Lat & Lon & Code & Group \\ \midrule
Cairns & 031011 & -16.87 & 145.75 & Am & Tropical \\
Darwin & 014015 & -12.42 & 130.89 & Aw & Tropical \\
Katherine & 014903 & -14.44 & 132.27 & Aw & Tropical \\
Townsville & 032040 & -19.25 & 146.77 & Aw & Tropical \\
AliceSprings & 015590 & -23.80 & 133.89 & BWh & Arid \\
Cobar & 048027 & -31.48 & 145.83 & BSh & Arid \\
Mildura & 076031 & -34.24 & 142.09 & BSk & Arid \\
Nhil & 078015 & -36.31 & 141.65 & BSk & Arid \\
SalmonGums & 012071 & -32.99 & 121.62 & BSk & Arid \\
Uluru & 015635 & -25.19 & 130.97 & BWh & Arid \\
Woomera & 016001 & -31.16 & 136.81 & BWh & Arid \\
Albury & 072160 & -36.07 & 146.95 & Cfa & Subtrop. \\
BadgerysCreek & 067108 & -33.90 & 150.73 & Cfa & Subtrop. \\
Brisbane & 040913 & -27.48 & 153.04 & Cfa & Subtrop. \\
CoffsHarbour & 059040 & -30.31 & 153.12 & Cfa & Subtrop. \\
GoldCoast & 040764 & -27.94 & 153.43 & Cfa & Subtrop. \\
Moree & 053115 & -29.49 & 149.85 & Cfa & Subtrop. \\
Newcastle & 061055 & -32.92 & 151.80 & Cfa & Subtrop. \\
NorahHead & 061366 & -33.28 & 151.58 & Cfa & Subtrop. \\
Penrith & 067113 & -33.72 & 150.68 & Cfa & Subtrop. \\
Richmond & 067105 & -33.60 & 150.78 & Cfa & Subtrop. \\
Sydney & 066062 & -33.86 & 151.21 & Cfa & Subtrop. \\
SydneyAirport & 066037 & -33.95 & 151.17 & Cfa & Subtrop. \\
WaggaWagga & 072150 & -35.16 & 147.46 & Cfa & Subtrop. \\
Williamtown & 061078 & -32.79 & 151.84 & Cfa & Subtrop. \\
\bottomrule\end{tabular}
\hfill
\begin{tabular}{llrrll}\toprule Location & BoM & Lat & Lon & Code & Group \\ \midrule
Wollongong & 068228 & -34.37 & 150.93 & Cfa$^\ast$ & Subtrop. \\
Adelaide & 023090 & -34.92 & 138.62 & Csb & Medit. \\
Albany & 009741 & -34.94 & 117.80 & Csb & Medit. \\
Dartmoor & 090194 & -37.92 & 141.26 & Csb & Medit. \\
MountGambier & 026021 & -37.75 & 140.77 & Csb & Medit. \\
Nuriootpa & 023373 & -34.48 & 139.01 & Csb & Medit. \\
PearceRAAF & 009053 & -31.67 & 116.02 & Csa & Medit. \\
Perth & 009225 & -31.92 & 115.87 & Csa & Medit. \\
PerthAirport & 009021 & -31.93 & 115.98 & Csa & Medit. \\
Portland & 090171 & -38.31 & 141.47 & Csb & Medit. \\
Walpole & 009998 & -34.95 & 116.72 & Csb & Medit. \\
Witchcliffe & 009746 & -34.03 & 115.10 & Csb & Medit. \\
Ballarat & 089002 & -37.51 & 143.79 & Cfb & Oceanic \\
Bendigo & 081123 & -36.74 & 144.33 & Cfb & Oceanic \\
Canberra & 070351 & -35.31 & 149.20 & Cfb & Oceanic \\
Hobart & 094029 & -42.89 & 147.33 & Cfb & Oceanic \\
Launceston & 091237 & -41.42 & 147.12 & Cfb & Oceanic \\
Melbourne & 086071 & -37.81 & 144.97 & Cfb & Oceanic \\
MelbourneAirport & 086282 & -37.67 & 144.83 & Cfb & Oceanic \\
MountGinini & 070349 & -35.53 & 148.77 & Cfb & Oceanic \\
Sale & 085072 & -38.12 & 147.13 & Cfb & Oceanic \\
Tuggeranong & 070339 & -35.42 & 149.09 & Cfb & Oceanic \\
Watsonia & 086068 & -37.74 & 145.10 & Cfb & Oceanic \\
NorfolkIsland & 200288 & -29.04 & 167.94 & Ocean & -- \\
\bottomrule\end{tabular}
\end{table}

\begin{table}[!htbp]
  \centering\scriptsize\setlength{\tabcolsep}{3pt}
  \caption{Model configurations used in Section 3 (read from the fitted scikit-learn pipelines). All models use the
  20 core features; imputation and scaling are fitted on the training rows only. Fit time: one fit on the random 70\%
  (about 100,000 days), 2 CPU cores.}\label{tab:hyper}
  \begin{tabular}{lp{4.3cm}p{6.6cm}r}
    \toprule
    Model & Preprocessing & Key hyperparameters & Fit (s) \\
    \midrule
    Logistic regression & median impute, standardise & multinomial, L2, $C=1$, lbfgs, max\_iter 3000 & 2.4 \\
    LR + polynomial & median impute, standardise, degree-2 polynomial (230 features), standardise & L2, $C=0.1$, lbfgs & 20.5 \\
    Decision tree (tuned) & median impute, standardise & Gini, max\_depth 16, min\_samples\_leaf 20 (grid search, grouped CV) & 0.9 \\
    $k$-NN & median impute, standardise & $k=15$, uniform weights, Euclidean & 0.2 \\
    Random forest & median impute, standardise & 200 trees, min\_samples\_leaf 5, max\_features $\sqrt{d}$, bootstrap & 14.5 \\
    MLP & median impute, standardise & hidden (64, 32), ReLU, Adam, $\alpha=10^{-3}$, 60 epochs & 10.9 \\
    HGB (default) & median impute & learning rate 0.1, max\_iter 100, 31 leaves, min\_samples\_leaf 20, no L2; early stopping on 10\% of the training rows & 4.2 \\
    \midrule
    HGB random search & \multicolumn{3}{p{11.9cm}}{16 configurations from: learning rate $\sim$ log-uniform(0.02, 0.3), max\_leaf\_nodes $\in[8,64)$, min\_samples\_leaf $\in[20,400)$, L2 $\sim$ log-uniform($10^{-3}$, 10), max\_iter $\in[100,400)$; station-grouped 5-fold CV on the 33 grouped-training stations.} \\
    \bottomrule
  \end{tabular}
\end{table}

\begin{figure}[!htbp]
  \centering
  \includegraphics[width=0.68\linewidth]{p3_k_selection}
  \caption{Day-level clustering for $k=2,\dots,10$: internal and external metrics, K-means inertia and GMM BIC.
  Subsampling stability (10 $\times$ 80\%): K-means 0.99, 0.98, 0.94 for $k=2,3,4$ and 0.69--0.89 for $k\ge5$;
  GMM 0.90 at $k=3$ and 5, otherwise 0.37--0.80; Ward 0.26--0.36 for every $k$ (notebook Section 3.2).}\label{fig:ksel}
\end{figure}

\end{document}
